\documentclass[11pt]{article}

\usepackage[margin=1in]{geometry}
\usepackage{amsmath, amssymb}
\usepackage{algorithm}
\usepackage{algpseudocode}
\usepackage{booktabs}
\usepackage{hyperref}

\title{Numerical Analysis I - Notes}
\author{}
\date{}

\begin{document}
\maketitle

\section{Numbers, Problems and algorithms}
 We are more interested in getting the relative error than the absolute error. The relative error is defined as
\[  
\frac{|x - x_{\text{true}}|}{|x_{\text{true}}|}
\]
where $x$ is the computed value and $x_{\text{true}}$ is the true value.

WE can use conditional numbers to determine how good our computation is. The conditional number is defined as
\[
\kappa = \frac{\|J\| \cdot \|x\|}{\|f(x)\|}
\]
where $J$ is the Jacobian matrix of the function $f$.

\subsection{Horner's algorithm}


\[
f(x) = a_0 + a_1 x + a_2 x^2 + \cdots + a_n x^n
\]
can be evaluated using Horner's algorithm:
\[
f(x) = a_0 + x(a_1 + x(a_2 + \cdots + x(a_{n-1} + x a_n) \cdots))
\]

\end{document}
